Released mathematics items carry a published tag, a standard code or a
cognitive domain, and a correct answer is read as evidence of that
competency, a reading that assumes passing requires it. We test the
requirement by searching for assessment bypasses: solver policies that
lack the tagged competency, by an explicit restriction on their
operations or information, yet pass under the published scoring rule. A
bypass is one-sided: it does not say students failed to learn, and
finding none does not certify the tag.

On released items from nine authorities we score three families, each
Holm-corrected: item-writing rules, test-taking strategies, and language
models shown the question with its given quantities withheld. Claimed
results are measured on items re-read from the printed page, which the PDF
text layer often misstates, and for the language models confirmed by a
second transcription. Item-writing rules yield no family-wise discovery;
one beats chance on two cells and awaits replication. Language models of
two families recover New York Regents Algebra~I and Algebra~II keys, and
the recovery survives a change of every option value with the form kept.
On items two independent codings mark as requiring computation on the
givens, one of three scorers misses the bar; the bypasses are
information-restricted, so a matched pass does not require computing on
the givens, and whether it requires the tagged operation is left open.
The international cognitive-domain cells, at about 50 to 80 public items
each, yield no bypass.
